\documentclass[10pt,a4paper]{amsart}
\usepackage[margin=19mm]{geometry}
\usepackage{amsmath,amssymb,mathtools}
\usepackage[scr=rsfs,cal=euler]{mathalfa}
\usepackage{xfrac}
\usepackage[unicode,hidelinks,bookmarks=false]{hyperref}
\newcommand{\ie}{i.e.\ }
\newcommand{\cf}{cf.\ }
\newcommand{\resp}{respectively\ }
\newcommand{\Subsection}{\subsection}
\providecommand{\nobreakdash}{\nobreak\hbox{-}}
\setlength{\emergencystretch}{2em}
\multlinegap0pt
\usepackage{amssymb}
\usepackage{mathtools}
\usepackage{tikz}
\usepackage{float}
\newtheorem{proposition}{Proposition}
\newcommand{\F}{\mathbb{F}}
\newcommand{\II}{\mathcal{I}}
\newcommand{\IM}{\mathrm{IM}}
\newcommand{\abs}[1]{\left|#1\right|}
\title{Large point-line matchings and small Nikodym sets}
\author{Zach Hunter, Cosmin Pohoata, Jacques Verstraete, Shengtong Zhang}
\date{Source: arXiv:2601.19879v1 (CC BY 4.0)}
\begin{document}
\maketitle

\noindent\emph{Source and licence.} Large point-line matchings and small Nikodym sets. Version arXiv:2601.19879v1 (CC BY 4.0), distributed by arXiv under the Creative Commons Attribution 4.0 International licence, http://creativecommons.org/licenses/by/4.0/, as recorded in the arXiv licence metadata for this exact version and shown on the abstract page https://arxiv.org/abs/2601.19879v1. Full licence notice: \texttt{licenses/arxiv-2601.19879-CC-BY-4.0.txt}.\par\smallskip

\noindent\textit{Editorial scope.} Counted unit: the article's proposition prop:dth-power-highd (source0.tex lines 543--553). The article's complete original proof of it is delivered with it (source0.tex lines 555--631, one \texttt{proof} environment of 77 lines). No mathematics is written, repaired or re-derived: the statement and the proof are the article's own lines, in the article's own notation, and no retained line is substituted or re-flowed.  No further source-local context is needed: every cross-reference of every retained line resolves inside the two delivered spans.  Nothing was omitted from any retained span, so the package carries no omission record.  No retained line contains a citation, so the wrapper carries no bibliography and no bibliography entry.  No retained line contains a figure environment, an image-inclusion command or drawing code, so no image is delivered and the package declares no asset dependency.  Editorial formatting only, in addition: an AMS \texttt{amsart} preamble that restates the article's own declarations and macros used by the retained lines, namely the environments \texttt{proposition} on separate counters, together with the standard packages those lines need.  The retained environments print in this document as Proposition 1 in order of appearance, whereas the article prints the same items with the numbering of its own full document.  The complete native source of the article as distributed in the arXiv e-print for this exact version is bundled unchanged as \texttt{sources/193495/source0.tex}.
\begin{proposition}\label{prop:dth-power-highd}
Assume $\mathrm{char}(\F_q)>d$.
Let $I\subset \F_q$ satisfy
\[
  (I-I)\cap D_d^\times=\emptyset.
\]
Then we have
\[
  \IM(d,q)\ \ge\ \frac{1}{(d-1)!}\,|I|\,(q-1)^{d-1}.
\]
\end{proposition}

\begin{proof}
Set
\[
  \Phi(x_1,\dots,x_d):=x_1+x_2^2+\cdots+x_d^d.
\]
Our goal is to construct a set of points $P\subset \F_q^d$ and, for each $p\in P$, a line $\ell_p$ such that $\ell_p\cap P=\{p\}$. The pairs $(p,\ell_p)$ would then form an induced matching in $\II_q^{(d)}$ of size $|P|$.

To do so, we will first build $(d-1)$-tuples $(x_2,\dots,x_d)\in(\F_q^\times)^{d-1}$ together with a direction vector
$v=(v_1,\dots,v_d)\in(\F_q^\times)^{d-1}\times\{1\}$
 (a vector $v$ which will be a function of $x_1,\ldots,x_{d}$) such that for all $x_1,\lambda \in \mathbb{F}_{q}$, we have
\begin{equation}\label{eq:Phi-shift}
  \Phi(x_1+\lambda v_1,\dots,x_d+\lambda v_d)=\Phi(x_1,\dots,x_d)+\lambda^d.
\end{equation}
Expanding by the binomial theorem, \eqref{eq:Phi-shift} is equivalent to requiring that all coefficients
of $\lambda^r$ for all $1\le r\le d-1$ vanish, while the $\lambda^d$ coefficient equals $1$ (the coefficient of $\lambda^d$ is $v_d^d=1$).

For $1\le r\le d-1$, the $\lambda^r$ coefficient equals
\[
  v_r^r\;+\;\sum_{i=r+1}^d \binom{i}{r}\,x_i^{\,i-r}v_i^r,
\]
so we want
\begin{equation}\label{eq:vr-rec}
  v_r^r = -\sum_{i=r+1}^d \binom{i}{r}\,x_i^{\,i-r} v_i^r.
\end{equation}

We choose the variables in decreasing order $r=d-1,d-2,\dots,2$, and handle $r=1$ at the end. When $r=d-1$, the sum in \eqref{eq:vr-rec} contains only the term $i=d$. This gives
$$v_{d-1}^{\,d-1} = -{d \choose d-1} x_d^{\,1}v_d^{\,d-1} = -d x_d,$$
since $v_{d}=1$. This provides the base case of the recursion.

For $2\le r\le d-2$, now let us assume that the variables $x_{r+2},\dots,x_d$ and $v_{r+1},\dots,v_d$ have already been chosen, and are nonzero. Then in \eqref{eq:vr-rec} the only occurrence of the new variable $x_{r+1}$ comes from the term with $i=r+1$. For the reader's convenience, we single out this term
\[
-\sum_{i=r+1}^d \binom{i}{r}\,x_i^{\,i-r} v_i^r
= -\binom{r+1}{r}\,x_{r+1}\,v_{r+1}^r \;-\; \sum_{i=r+2}^d \binom{i}{r}\,x_i^{\,i-r} v_i^r.
\]
This allows us to emphasize that the right-hand side of \eqref{eq:vr-rec} is an affine-linear function of $x_{r+1}$ of the form
\[
A\,x_{r+1}+B,\ \ \ \ \ \text{where}
\qquad
A:=-(r+1)v_{r+1}^r,\quad
B:=-\sum_{i=r+2}^d \binom{i}{r}\,x_i^{\,i-r} v_i^r.
\]
Note that $B$ is already determined by the previously chosen variables.
Moreover $A\neq 0$ because $v_{r+1}\neq 0$ and $\mathrm{char}(\F_q)>d\ge r+1$.
Therefore the map $x_{r+1}\mapsto A x_{r+1}+B$ is a bijection of $\F_q$, so as $x_{r+1}$ ranges over $\F_q$
the right-hand side of \eqref{eq:vr-rec} ranges over all of $\F_q$.

In particular, we may choose $x_{r+1}\in\F_q^{\times}$ so that the right-hand side lies in $D_{r}^{\times}$ (recall $D_{r}^{\times}=\left\{t^r:t\in\F_q^\times\right\}$),
and then pick $v_r\in\F_q^\times$ satisfying \eqref{eq:vr-rec}.
Note that
$$\abs{D_r^{\times}} = \frac{q-1}{\gcd(r,q-1)} \geq \frac{q - 1}{r},$$
giving at least $(q-1)/r$ valid choices of $x_{r+1}$ for which the right-hand side
of \eqref{eq:vr-rec} lies in $D_{r}^{\times}$. For each of these, we can then choose $v_r\in\F_q^\times$ satisfying \eqref{eq:vr-rec}. 

After choosing $v_2,\dots,v_d$, we choose $v_1\in\F_q$ so that the $\lambda^1$ coefficient vanishes;
this is always possible because the $r=1$ equation \eqref{eq:vr-rec} is linear in $v_1$.

We have thus obtained a set $S\subset(\F_q^\times)^{d-1}$ of admissible tuples $(x_2,\dots,x_d)$, and for each $(x_2,\dots,x_d)\in S$ we fixed one corresponding direction $v(x_2,\dots,x_d)$ satisfying \eqref{eq:Phi-shift}. Moreover, we have also shown along the way that
$$|S|\ \ge\ \prod_{r=1}^{d-1}\frac{q-1}{r} = \frac{(q-1)^{d-1}}{(d-1)!}.$$
We can define finally the point set
$$P = \left\{p=(x_1, x_2, \cdots, x_d) \in \mathbb{F}_{q}^{d}:\ (x_2, \cdots, x_t) \in S\ \ \text{and}\ \ \Phi(p) = x_1 + \sum_{i=2}^d x_i^i \in I\right\}.$$
By construction, we have
\[
  |P|=|I|\,|S|\ \ge\ \frac{1}{(d-1)!}\,|I|\,(q-1)^{d-1}.
\]
For $p = (x_1, x_2, \cdots, x_t) \in P$ arising from $(x_2,\dots,x_d)\in S$, define the affine line
\[
  \ell_p:=\{p+\lambda v(x_2,\dots,x_d):\ \lambda\in\F_q\}.
\]
By \eqref{eq:Phi-shift}, for every $\lambda\in\F_q$ we have
$\Phi(p+\lambda v)=\Phi(p)+\lambda^d$.
If $p'\in \ell_p\cap P$, then $\Phi(p'),\Phi(p)\in I$ and
\[
  \Phi(p')-\Phi(p)=\lambda^d\in D_d.
\]
By the hypothesis $(I-I)\cap D_d^\times=\emptyset$, this forces $\lambda=0$, hence $p'=p$.
Thus $\ell_p\cap P=\{p\}$ for every $p\in P$, and $\{(p,\ell_p)\}_{p\in P}$ is an induced matching of size $|P|$.
\end{proof}

\end{document}
